% !TeX root = ../tg.tex

\chapter{研究方法与技术路线}
\section{引言}
本章旨在立足于分布式学习的研究范式，系统阐述本文所采用的研究方法及技术路线的创新之处。内容主要涵盖两个维度：首先，详述分布式研究中核心的理论证明方法与实验验证体系；其次，从技术演进视角出发，深入剖析本文如何在继承现有研究脉络的基础上实现突破，从而构建出兼具延续性与创新性的技术路径。
% \section{研究方法}
\section{理论证明}
本章重点聚焦于分布式学习的收敛性理论证明。正如~\ref{subsec:bp-convergence}节所述，当前分布式学习领域已涌现出大量基于联邦学习框架的收敛性分析方法。
在此基础上，本文沿袭了 Reddi 等人~\cite{reddi2020adaptive}的经典证明范式，通过设定梯度期望平方、梯度方差以及数据异构性的上界约束，深入探究在Non-IID数据与标签噪声干扰下，上述上界的演化规律。
基于此分析，本文进一步推导并确立了相应复杂场景下SFL的收敛性上界。
\section{实验探究}
\subsection{实验方法}
遵循分布式学习领域的通用研究范式，本文的实验体系由仿真实验与测试平台实验两部分构成。
仿真实验旨在通过软件模拟构建 SFL 的客户端与服务器组件，在可控的虚拟环境中验证算法各模块的协同机制与逻辑正确性；其核心评估指标聚焦于模型的准确率与损失。
测试平台实验则基于真实硬件设备搭建分布式测试平台，将算法中的各个逻辑组件映射至独立的物理节点，以在真实的网络与环境约束下评估系统性能；其关键指标涵盖能耗、计算延迟、内存占用及通信开销等资源维度。
在实施条件上，仿真实验通常仅需高性能服务器即可完成；而测试平台实验则需依赖树莓派集群及配套网络设备，以模拟真实的分布式计算场景。

\subsection{数据集}
遵循分布式学习领域的通用实验范式，本文选取了 CIFAR-10、CIFAR-100 及 Fashion-MNIST (FMNIST) 三个基准分类数据集进行评估。
在数据特性方面，CIFAR-10 与 CIFAR-100 均为三通道彩色图像数据集，各包含 60,000 张样本（其中训练集 50,000 张，测试集 10,000 张），类别数分别为 10 类和 100 类；FMNIST 则为单通道灰度图像数据集，共包含 [填入总数，如 70,000] 张样本（训练集 [填入数量] 张，测试集 [填入数量] 张），涵盖 10 个服饰类别。
针对实验场景的适配性，本文采取了差异化部署策略：由于 CIFAR 系列数据集具有高分辨率、高复杂度及大容量的特征，其推理过程对计算资源要求较高，超出了树莓派等边缘设备的承载能力，因此仅用于仿真实验以验证算法的理论性能。相比之下，FMNIST 作为一种轻量级数据集，不仅兼容边缘设备的资源约束，适用于测试平台实验；同时，相较于传统的 MNIST 手写数字数据集，FMNIST 具有更高的分类难度，能够更有效地反映模型在复杂特征提取下的准确率与损失收敛情况，因而也被纳入仿真实验的评估体系中。

\subsection{低质量数据的构造}
\subsubsection{Non-IID数据}
遵循分布式学习领域的通用实验范式，本文采用\textbf{狄利克雷分布（Dirichlet Distribution）}来模拟客户端间数据的非独立同分布（Non-IID）特性。在该分布模型中，浓度参数 $\alpha$ 起着关键调控作用：$\alpha$ 值越小，生成的样本标签分布向量越稀疏，意味着不同用户间的数据分布差异越大（即异构性越强）；反之，$\alpha$ 值越大，数据分布则趋向于均匀（IID）。

具体而言，假设数据集中共有 $K$ 个类别，对于第 $n$ 个用户（Client $n$），其本地数据集的类别比例分布向量 $\mathbf{p}_n = [p_{n,1}, p_{n,2}, \dots, p_{n,K}]^\top$ 服从参数为 $\alpha$ 的狄利克雷分布：
\begin{equation}
    \mathbf{p}_n \sim \text{Dir}(\alpha \cdot \mathbf{1}_K)
\end{equation}
其中，$\mathbf{1}_K$ 表示长度为 $K$ 的全 1 向量。该分布的概率密度函数定义为：
\begin{equation}
    f(\mathbf{p}_n; \alpha) = \frac{1}{B(\alpha)} \prod_{i=1}^{K} p_{n,i}^{\alpha - 1}
\end{equation}
式中，$B(\alpha)$ 为归一化常数（多元 Beta 函数），且满足约束条件 $\sum_{i=1}^{K} p_{n,i} = 1$ 及 $p_{n,i} \geq 0$。

在实际数据划分过程中，我们首先针对每个用户独立采样得到其类别比例向量 $\mathbf{p}_n$，随后依据该比例将全局数据集中的样本分配至各用户。通过调整 $\alpha$ 的取值（例如 $\alpha \in \{0.1, 0.5, 1.0\}$），我们可以灵活构建从极度 Non-IID 到近似 IID 的不同实验场景，从而全面评估算法在各类数据异构程度下的鲁棒性。

\subsubsection{标签噪声}
我们在实验中使用了两种方法构造标签噪声，一种是随机翻转，另一种是类别相关翻转。
